\section{Discussion}\label{sec:discussion}

\paragraph{Why MADII falls short} There are three reasons, and none of them requires a bug in our rebuild.
\begin{itemize}[leftmargin=*]
  \item \textbf{The imitation stage does not teach the expert's behaviour.} It only fills the replay memory with expert decisions for Q-learning, and the network is never trained to reproduce them. Consistent with this, a run that used battery Dijkstra (a 151-round expert) as the imitation source fell from 7.6 to 4.4 rounds between episodes 25 and 50; this run was stopped early.
  \item \textbf{The Informer has nothing to forecast.} It was designed for long time-series forecasting~\cite{zhou2021informer}. MADII feeds it a single snapshot of 101 nodes and asks it to route, and its sparse attention gains nothing at that size.
  \item \textbf{No baseline revealed the gap.} MADII is compared only with clustering, metaheuristic and learned methods, so the distance to shortest-path routing and to the optimum went unnoticed.
\end{itemize}
The broader lesson concerns evaluation practice. A learned router should be reported against a battery-weighted shortest-path router and against a provable ceiling. Both are cheap to compute, and together they place any new method on an absolute scale.

\paragraph{Why coordination, and why exact energy} Battery-weighted Dijkstra already captures most of the lifetime available from a single tree per round. What it lacks is coordination: nodes choose independently and overload shared relays. Booking load within the round supplies that coordination without a solver. Each extra unit of booked traffic raises a relay's cost smoothly, so close decisions tip towards alternatives while necessary relays are still used. The ablation shows that the two signals correct each other. Energy without load lets nodes pile onto one relay, and load without energy keeps relaying through nearly empty nodes. Energy must stay exact for a different reason: it drains slowly, so any reporting rule that is robust to noise, such as LEST's hysteresis trigger, freezes it. Load jumps from round to round and survives compression into 2 bits. The same freeze explains an observation from our earlier fuzzy-Q project, where every node remained in the same energy tier for the whole run.

\paragraph{When to use which router} If a sink can afford one LP solve per round (about 0.2~s at $n=100$), the LP planner is the best choice, at 97\% of the oracle. If it cannot, LEST Dijkstra reaches about 90\% with shortest paths alone, and battery Dijkstra about 85\% with a single Dijkstra run. A traffic forecast is a small robustness refinement for the planner and does not help Dijkstra. In this setting, learning brings no benefit we could measure.

\subsection{Threats to validity}\label{sec:threats}
\paragraph{Internal validity} MADII is our rebuild from the paper's text. It reproduces the headline lifetime and delivered data, but it is not the authors' code, which has not been released; we have requested it. Two harnesses give MADII different FND values (19.5 and 27.3 rounds); both are below 17\% of the ceiling, and the classical routers agree across harnesses. EBR-DA is evaluated outside the aggregation setting it was designed for, with buffer occupancy approximated by load. All tuned settings were chosen on validation deployments, and the test sets were run once.
\paragraph{Construct validity} FND is one of several lifetime definitions. We also report 10\%, 25\% and 50\% dead, delivered data and energy efficiency (Table~\ref{tab:traffic}), and the ranking is the same.
\paragraph{External validity} The radio model is the idealised one used by MADII (A3--A6): no collisions, no lossy links, no barriers and no idle-listening energy. Absolute percentages may shift under a packet-level MAC with retransmissions. Lossy links would raise the cost of long hops and of busy relays alike, which we expect to favour load-aware routing, but this has not been tested. The synthetic traffic models are designed by us. The Intel layout is scaled, its traffic is derived from readings rather than measured, and test deployments overlap in time. Networks of 54 and 100 nodes are tested; the LP grows with the square of $n$, and LEST Dijkstra's worst case with its cube. Both LEST Dijkstra and the LP planner compute routes at the sink. A fully distributed variant, in which every node computes the same tree from a shared snapshot, would also need every battery level at every node, and it has not been tested.
